% =====================================================================
%  Chapter 20 — Five Senses of Tapping
% =====================================================================

\chapter{Five Senses of Tapping}
\label{ch:senses}

The middle clause of this book's title asks whether we have already tapped
dark matter. The question is not a riddle, and it is not rhetorical: it has
five distinct answers depending on what is meant by the word, and the
answers differ in ways that matter for what the field should do next.

\section{The five senses}The framework used in this book distinguishes the
following.

\begin{itemize}
\item \textbf{Detection through gravity} --- using dark matter as a probe of
its own gravitational effect on light and on the motions of ordinary matter.
This has been done since 1933 and is done to high precision.
\item \textbf{Detection through a local enhancement} --- using the
gravitationally captured population in the Solar System, where the density
is four to five orders of magnitude above the galactic value, and where the
velocity distribution is different from the halo's.
\item \textbf{Conversion into energy} --- using the dark matter as a source
of usable power, by capturing it and letting it annihilate or decay, or by
converting it electromagnetically.
\item \textbf{Conversion into a signal in an engineered device} --- using
the ambient field to produce a measurable output in a quantum system.
\item \textbf{Detection as a particle} --- producing an event in a detector
that can only be attributed to a discrete dark matter particle.
\end{itemize}

\begin{table}[H]
\centering
\caption{The five senses of tapping, with the status of each as of November
2026.}
\label{tab:senses}
\begin{tabular}{@{}>{\RaggedRight\arraybackslash}p{0.30\textwidth}p{0.17\textwidth}p{0.39\textwidth}@{}}
\toprule
\textbf{Sense} & \textbf{Status} & \textbf{The qualification} \\
\midrule
Gravity & Yes & Done since 1933, at high precision; a tool rather than a discovery \\
Capture & Yes & A four-to-five order enhancement, in principle directly measurable \\
Harvesting & Physically allowed & The coupling is the obstacle; the arithmetic is absurd \\
Quantum devices & Partly & A predicted power in a qubit; far too small to be useful as energy \\
Particle & No & No dark matter particle has ever been observed \\
\bottomrule
\end{tabular}
\end{table}

\tabref{tab:senses} is the short answer to the title question, and the
qualification column is where the interesting content sits. Three of the
five senses have been achieved. The two that have not --- harvesting as
energy, and the particle --- are separated by the same quantity: the size of
the interaction.

\section{Why the distinction is worth making}A single answer to the question
would either overstate the field's achievement or dismiss it.

If the question means ``has anyone detected dark matter'', the answer is
that the detection has been gravitational since 1933, the evidence is
overwhelming, and the particle has never been seen. If it means ``is there
any sense in which we have made use of dark matter', the answer is yes: it
is the instrument with which the cosmic microwave background was measured
and the cosmic web was mapped, and the density it implies is known to two
per cent.

The five senses also order the field's priorities, because each sense has a
different bottleneck. Gravity has no bottleneck; it is finished, and it
constrains the total abundance and the large-scale distribution but nothing
about the particle. Capture has a bottleneck in detector technology and in
the difficulty of distinguishing the bound population from the unbound one.
Harvesting has a bottleneck that is a pure number, and that number is
roughly ten orders of magnitude short of what a power plant needs.
Quantum devices have a bottleneck in signal-to-noise. The particle has a
bottleneck that is the neutrino fog of Chapter~\ref{ch:direct}.

\section{The uncomfortable structure}Two facts sit awkwardly together and the rest
of Part~V is devoted to them.

The first is that the senses with positive answers are the ones that do not
require the particle. The second is that the one sense with a negative
answer is the one that everything else depends on. Without the particle,
there is no way to know what the dark matter is, no way to confirm any of
the candidates, and no way to decide what to build next.

That is the actual content of the title's third clause, and Chapter~\ref{ch:roadmap}
is where it becomes a programme.

\begin{keytakeaways}
\item Five senses of tapping have distinct answers: gravity yes, capture
yes, harvesting physically allowed but absurd in practice, quantum devices
partly, the particle no.
\item The senses that have succeeded are the ones that do not require the
particle, and the one that has not is the one everything else depends on.
\item Each sense has a different bottleneck, and the bottlenecks determine
where the field's effort should go.
\end{keytakeaways}
